% Table (F7) — what each family CAN (green check) and CANNOT (red cross) do.
\begin{table*}[t]
\centering
\footnotesize
\renewcommand{\arraystretch}{1.3}
\begin{tabular}{@{}p{0.19\textwidth} *{6}{c} p{0.23\textwidth}@{}}
\toprule
\textbf{Family} & \textbf{Gen.} & \textbf{Inspect.} & \textbf{Adapt.} & \textbf{Compos.}
 & \textbf{Distrib.} & \textbf{Self-impr.} & \textbf{Key limitation} \\
\midrule
End-to-end VLA (\S\ref{sec:vla}) & \yy & \nn & \nn & \nn & \pp & \nn & no inspectable reuse or editing \\
Code-as-policy (\S\ref{sec:code}) & \pp & \yy & \yy & \yy & \pp & \yy & brittle under perception errors \\
Reward synthesis (\S\ref{sec:reward}) & \pp & \nn & \yy & \nn & \nn & \pp & ships opaque weights, sim-bound \\
RL skill discovery (\S\ref{sec:skill}) & \pp & \nn & \pp & \pp & \nn & \nn & symbolically un-inspectable \\
Market skills (\S\ref{sec:open}) & \nn & \nn & \nn & \pp & \yy & \nn & no adaptation or provenance \\
\bottomrule
\end{tabular}
\caption{What each family can (\yy) and cannot (\nn) do (\pp\ partial). Columns:
\textbf{Gen.}eralization, \textbf{Inspect.}ability, \textbf{Adapt.}ability,
\textbf{Compos.}ability, \textbf{Distrib.}utability, and run-time \textbf{Self-impr.}ovement.
No family scores well on all six, the tension the skill economy (\S\ref{sec:open}) inherits.}
\label{tab:branch_limits}
\end{table*}
